\section{Related Work}
\label{sec:related work}

\subsection{Conversation and Speech Video Understanding}
The task of understanding dialogue and speeches in videos has always been a widely explored issue~\cite{han2021bi,liang2021multibench,zhang2021multi,4130374,10147343}. Most of the datasets~\cite{hasan2019ur,poria2018meld} and methods in this field primarily focus on emotion analysis for specific tasks, identifying the emotions, humor, or sarcasm of the speaker from dialogue, speech videos, and subtitle texts. It has been indicated by Chauhan et al.~\cite{chauhan2020sentiment} that there is a close connection between sentiment analysis~\cite{li2020multi,song2022supervised}, sarcasm detection~\cite{castro-etal-2019-towards}, and humor detection~\cite{islam2021multi}.

Many methods~\cite{chauhan2020sentiment,ray2022multimodal,hazarika2020misa} utilized pre-trained Convolutional Neural Network~(CNN)~\cite{li2021survey} and Bidirectional Encoder Representations from Transformers (Bert)~\cite{devlin2018bert} to extract textual and video features, followed by the use of a Multi-Modal Transformer~\cite{tsai2019multimodal} for feature fusion. However, these methods required training attention models from scratch, and it has been shown that attention mechanisms perform poorly with limited data. Additionally, these models are hindered by a lack of domain-specific and general knowledge, failing to fully understand advanced linguistic expertise and visual information. With the development of LLMs achieving significant success in the NLP field, InstructERC~\cite{lei2023instructerc} proposed a method using LLMs to address emotion recognition tasks. By leveraging the knowledge of large models, it achieved commendable results, but it only considered the textual modality and did not utilize the video modality. Furthermore, it required finetuning the LLM, which incurs high training costs.

%Video sarcasm and humor detection are two closely related sentiments, transitioning from text-centric to multimodal~\cite{cheng2021multimodal, 10192323} methods, represents a significant stride in understanding complex human sentiments. Due to transient nature and data scarcity, the challenges of identifying sarcasm and humor in visual media  highlight the need for multimodal systems~\cite{pramanick2022multimodal}. These systems effectively harness self-attention and optimal transport to navigate intra- and cross-modal nuances. The nuanced expression of sarcasm and humor, often conveyed through a blend of verbal and non-verbal cues, underscores the intricacy of these sentiments \cite{bedi2021multi}. Some work~\cite{dong2023simmmdg,li2023boosting} uses multimodal fusion methods in sarcasm detection. Pioneering datasets like UR-FUNNY~\cite{hasan-etal-2019-ur} and MUStARD \cite{mustard} have been instrumental in advancing multimodal humor and sarcasm detection.

\subsection{Video Large Language Models}

Large language models such as ChatGPT~\cite{liu2023summary} and LLAMA~\cite{touvron2023llama} have achieved tremendous success in the NLP field. Many works (like LLAVA~\cite{liu2023visual} and GPT-4~\cite{liu2023summary}) attempt to integrate other modalities such as vision and audio to enhance understanding~\cite{chai2022deep,wang2023user}. LLaVA uses adapters~\cite{he2021effectiveness} to map image information tokens directly into the natural language representation space. BLIP-2~\cite{li2023blip}, leveraging a frozen LLM and Image Encoder, introduces a feature extraction module Q-former for multi-modal fusion, based on text-image contrastive learning. This significantly enhances the capability of multi-modal feature representation. Building on the success of image and LLM integration, many works~\cite{song2023moviechat,maaz2023video,zhang2023video} have started focusing on applying LLMs to video understanding, especially in video question answering (VQA) tasks (like VideoLLAMA~\cite{zhang2023video}). VideoLLAMA introduced a temporal fusion Q-former for modeling temporal information. VideoChat~\cite{maaz2023video} uses a Video Foundation Model~\cite{arnab2021vivit} to obtain video representations and align them with LLMs. However, all these models underwent pre-training on extensive video datasets, making the training costs unaffordable for most researchers and practitioners. Our proposed MTransLLAMA obviates the need for video pre-training and costly temporal modules, achieving a transition from an image-based LMM to a video-based LMM. Moreover, these models' video pre-training datasets are primarily VQA datasets, which are not efficient for other downstream video understanding tasks like emotion analysis with a large domain gap.

% Vision-Language Models (VLMs) has showcased remarkable progress, proving highly effective in addressing a myriad of challenges with superior performance \cite{visualgpt,kosmos,yang2022zero,tiong2022plug,alayrac2022flamingo,blip2,blip1,palm_e}. Models such as Flamingo \cite{alayrac2022flamingo} and BLIP-2 \cite{blip2} exemplify the strides made in sophisticated image-text alignment and contextual understanding. The advent of GPT-4 \cite{gpt4} and PaLM-E \cite{palm_e} further demonstrates the expansive capabilities of VLMs in visual reasoning and real-world sensor integration. Despite these advancements, the application of VLMs in video sarcasm detection remains largely unexplored, presenting a significant opportunity. Leveraging VLMs in this domain could provide breakthroughs in detecting nuanced expressions and contextual incongruities in videos, a frontier yet to be fully tapped.

\begin{figure*}[t]
    \centering
    \includegraphics[width=\textwidth]{figs/framework2.pdf}
    \caption{\textbf{Overview of MTransLLAMA}. 
    % When the input text conversation length is $T$ for text data $C$, and the padding length is $K_T$. 
    % Firstly, we conduct spatial feature extraction with $K_Q$ pre-defined queries. After Spatial-Self Attention, we swap the channels of queries, entering Temporal Self Attention, where temporal information interaction occurs between tokens with similar semantics across different frames. Following Temporal Self Attention, we further transform the query channels and perform image feature extraction on the input $T$ frames of the video. Finally, we aggregate the last layer tokens of the Q-former, concatenate them with instructions, and input them into the frozen LLM to obtain the final natural language outputs. We introduce LoRA separately in different attention modules for fine-tuning while keeping all other parameters frozen. Additionally, we incorporate the DRT module for dynamic routing of attention scope within the attention modules.
    First, we perform spatial feature extraction using $K_Q$ pre-defined queries. After applying Spatial Self-Attention, we swap the channels of the queries and proceed to Temporal Self-Attention, where temporal information interaction occurs between tokens with similar semantics across different frames. Following Temporal Self-Attention, we further transform the query channels and extract image features from the input $T$ frames of the video. Finally, we aggregate the tokens from the last layer of the Q-former, concatenate them with instructions, and input them into the frozen LLM to generate the final natural language outputs. We introduce LoRA separately in different attention modules for fine-tuning, while keeping all other parameters frozen. Additionally, we incorporate the DAR module for dynamic routing of attention scope within the attention modules.}
    \label{fig:framework}
\end{figure*}


\subsection{Temporal Modeling}
With the advancement of AI, more attention is shifting from image to video understanding, including tasks like video question answering~\cite{li2022invariant,xiao2022video,10334011} and video action recognition~\cite{feichtenhofer2017spatiotemporal,wang2021actionclip,qi2022weakly}. Unlike image models, due to the temporal nature of videos, models require temporal modules for video understanding. With the advent of the Transformer, more work has utilized attention modules as the backbone for video understanding. With the rise of Image pre-trained models, ViT, CLIP~\cite{radford2021learning}, and their variants (like Deter~\cite{carion2020end} and DINO~\cite{caron2021emerging}) have been introduced, achieving state-of-the-art performance in image classification~\cite{deng2009imagenet,9917532,10283961}, object detection~\cite{zhao2019object,9786052,9548849}, and segmentation~\cite{minaee2021image}. The visual patterns they learn from large-scale image-text pre-training data are of great value.
However, due to the quadratic complexity growth of Transformers, large temporal convolution modules have introduced significant training costs for video understanding.

To overcome these limitations, a strategy known as parameter-efficient transfer~\cite{hu2021lora,he2021effectiveness} has become increasingly popular in NLP, aiming to fine-tune only a few parameters while keeping large pre-trained models frozen to achieve strong performance. As large ViTs develop, these techniques are being introduced into Computer Vision. However, these works~\cite{liu2023visual,bao2021beit} either focus on tuning pre-trained image models for image task finetuning or adjusting pre-trained video models for downstream video tasks. Consequently, many works have transferred image-based pre-trained models to the video domain. Some methods~\cite{fang2021clip2video} focus on prompt or sampling modeling, while other methods~\cite{luo2022clip4clip} design temporal modules as intermediate structures, as illustrated in the middle. Yang et al. (AIM)~\cite{yang2023aim} repurpose CLIP’s attention modules through adapters for efficient video transfer. However, these methods do not explore multiple modalities, and the temporal attention of image patches lacks interpretability. We propose Multi-Modal Query Fine-tuning, achieving modeling with a multi-modal Q-former without temporal modules.

% Parameter-efficient finetuning, initially introduced in natural language processing (NLP) \cite{adapter,hu2022lora,prompttuning,li-liang-2021-prefix,unifiedadapter,ben-zaken-etal-2022-bitfit,fixsparsemask,qing2022mar}, have emerged in response to the impracticality of fully finetuning increasingly larger language models for diverse downstream tasks. The primary aim of these methods is to minimize the number of trainable parameters, thereby reducing computational costs while maintaining or even enhancing the performance achieved through complete finetuning. In a similar vein, the realm of computer vision has begun exploring parameter-efficient transfer learning \cite{vpt,visualpixelprompt,chen2022adaptformer,convadapter,gao2022dept}, primarily focusing on model adaptations within the same domain, such as image-to-image or video-to-video transformations. Notably, AIM~\cite{yang2023aim} has innovated in this area by adapting pre-trained image transformer models for efficient video processing. Our approach advances this concept by adapting the image model to accommodate multimodal tasks.




